%% ============================================================
%% 3. METHOD — CRR-Flow: State-Aware Flow Matching for HOI
%% ============================================================
%% Replace the existing Section 3 (Method) in sample-sigconf-authordraft.tex
%% with this content.

\section{Method}
\label{sec:method}

Given a natural language instruction $P$ describing a manipulation task (potentially multi-sentence), our goal is to generate a temporally coherent human-object interaction motion $\mathcal{M} \in \mathbb{R}^{T \times D}$ depicting a full-body human manipulating one or more objects. Our framework, \textbf{CRR-Flow} (Contact-Release-Reach State-Aware Flow), consists of three stages: (1)~an LLM-based \emph{State Planner} that translates text into per-entity physical state timelines (\cref{sec:state_planner}), (2)~an \emph{Entity-Centric Diffusion Transformer} that generates motions conditioned on states, text, and object geometry via flow matching (\cref{sec:entity_dit}), and (3)~a \emph{Block-wise Autoregressive Relay} that extends generation to arbitrary horizons through overlapping kinematic history windows (\cref{sec:autoreg}).


\subsection{State Planning: Text to Physical State Maps}
\label{sec:state_planner}

We observe that human-object manipulation can be decomposed into a finite set of \emph{physical interaction states}: a hand is either idle (\textsc{Static}), approaching the target (\textsc{Reach}), in contact (\textsc{Grasped}), performing a tool action (\textsc{Interacting}), or withdrawing (\textsc{Release}). Objects correspondingly transition between \textsc{Static}, \textsc{Grasped}, and \textsc{Interacting}, while unused object slots are marked \textsc{Padding}.

\paragraph{Per-entity state labels.}
For each frame $t$, we define a state vector $\mathbf{s}_t \in \{0,1,2,3,4,5\}^7$ across 7 entity slots: body, left hand, right hand, and up to 4 objects. This yields a state timeline $\mathbf{S} \in \mathbb{Z}^{T \times 7}$ that serves as a ``choreography sheet'' specifying \emph{when} each entity should act, decoupling temporal structure from the continuous motion generation.

\paragraph{Ground-truth state extraction.}
During training, state labels are computed automatically from motion capture data using a Schmitt trigger contact detector with hysteresis:
\begin{equation}
    c_t = \begin{cases}
        \text{true} & \text{if } d_t < \delta_{\text{on}} \text{ (trigger)} \\
        \text{false} & \text{if } d_t > \delta_{\text{off}} \text{ (release)} \\
        c_{t-1} & \text{otherwise (deadzone)}
    \end{cases}
    \label{eq:schmitt}
\end{equation}
where $d_t$ is the minimum fingertip-to-object distance at frame $t$, $\delta_{\text{on}} = 1.5$~cm, and $\delta_{\text{off}} = 3.0$~cm. The deadzone between thresholds prevents contact chatter at boundaries. A majority-vote filter (window size 5) further smooths isolated state transitions.

Hand-object binding is determined from task annotations: the primitive's \texttt{obj\_list\_lh} and \texttt{obj\_list\_rh} fields specify which hand manipulates which object, while the action topology (triadic vs.\ dyadic) determines whether secondary objects receive \textsc{Interacting} labels.

\paragraph{LLM-based state inference.}
At test time, when ground-truth states are unavailable, a large language model (e.g., GPT-4) parses the text instruction into structured state timelines. Given the instruction and entity list, the LLM reasons about action phases and outputs a JSON state schedule:
\begin{equation}
    \text{LLM}(P, \text{entities}) \rightarrow \mathbf{S} \in \mathbb{Z}^{T \times 7}.
    \label{eq:llm_planner}
\end{equation}
This enables open-vocabulary generation without requiring ground-truth annotations.


\subsection{Entity-Centric Flow Matching DiT}
\label{sec:entity_dit}

\paragraph{Motion representation.}
We represent the human body using \textbf{77 surface markers} extracted from SMPL-X~\cite{pavlakos2019expressive} mesh vertices, following the SSM67 protocol~\cite{interact2024} augmented with 10 fingertip markers. The markers are partitioned into three groups: body (50 markers, 150D), left hand (14 markers, 42D), and right hand (13 markers, 39D), totaling 231D of positional features.

Each object is represented by a 9D pose vector (3D translation + 6D continuous rotation~\cite{zhou2019continuity}), with up to 4 object slots. Object geometry is encoded using \textbf{Basis Point Sets (BPS)}~\cite{prokudin2019efficient}: 1024 Fibonacci-distributed basis points, each storing the scalar distance to the nearest object surface point, yielding a static 1024D shape embedding per object.

\paragraph{Velocity features.}
For each positional feature, we compute the first-order temporal difference $\mathbf{v}_t = \mathbf{x}_t - \mathbf{x}_{t-1}$ \emph{after} temporal downsampling to the target frame rate. This produces a 267D velocity vector with the same layout as positions. The full dynamic feature is the concatenation $[\mathbf{x}_{\text{pos}}; \mathbf{x}_{\text{vel}}] \in \mathbb{R}^{534}$ per frame.

\paragraph{Entity tokenization.}
Unlike prior methods that treat the full feature vector as a monolithic token, we decompose it into \textbf{7 entity tokens} per frame, each processed through a dedicated input projection:
\begin{align}
    \mathbf{h}_{\text{body}} &= W_{\text{body}} \, \mathbf{x}_{\text{body}} \in \mathbb{R}^D, \quad (300 \rightarrow 512) \\
    \mathbf{h}_{\text{lh}} &= W_{\text{lh}} \, \mathbf{x}_{\text{lh}} \in \mathbb{R}^D, \quad (84 \rightarrow 512) \\
    \mathbf{h}_{\text{rh}} &= W_{\text{rh}} \, \mathbf{x}_{\text{rh}} \in \mathbb{R}^D, \quad (78 \rightarrow 512) \\
    \mathbf{h}_{\text{obj}_i} &= W_{\text{obj}} \, [\mathbf{x}_{\text{obj}_i}; \mathbf{b}_i] \in \mathbb{R}^D, \quad (1042 \rightarrow 512)
    \label{eq:entity_proj}
\end{align}
where $\mathbf{b}_i \in \mathbb{R}^{1024}$ is the \emph{clean} (non-noised) BPS embedding concatenated to the object's dynamic features before projection. The object projection $W_{\text{obj}}$ is shared across all 4 slots.

\paragraph{Per-entity state injection.}
State labels are embedded and added \emph{directly} to the corresponding entity token, not broadcast globally:
\begin{equation}
    \mathbf{h}_e \leftarrow \mathbf{h}_e + \text{StateEmbed}(\mathbf{s}_{t,e}), \quad e \in \{\text{lh}, \text{rh}, \text{obj}_1, \ldots, \text{obj}_4\}.
    \label{eq:state_inject}
\end{equation}
The body entity receives no state injection (always \textsc{Static}). This per-entity design resolves the ambiguity of global conditioning: the model knows exactly \emph{which} hand is reaching and \emph{which} object is being grasped at each frame.

\paragraph{Spatio-temporal self-attention.}
The 7 entity tokens per frame are stacked with dual positional encodings---sinusoidal temporal PE and learned entity PE---then flattened to a sequence of $T \times 7$ tokens:
\begin{equation}
    \mathbf{H} = \text{Flatten}\left(\mathbf{h}_{t,e} + \text{PE}_{\text{time}}(t) + \text{PE}_{\text{entity}}(e)\right) \in \mathbb{R}^{(T \cdot 7) \times D}.
    \label{eq:flatten}
\end{equation}
This sequence is processed by $L{=}8$ DiT blocks~\cite{peebles2023scalable} with AdaLN-Zero~\cite{peebles2023scalable} modulation and FlashAttention-2~\cite{dao2023flashattention2}, enabling efficient cross-entity attention: the left hand at frame $t$ can attend to the object at frame $t{\pm}k$, learning spatio-temporal coordination patterns.

Unused object slots (state $= \textsc{Padding}$) are masked via key padding, preventing the model from hallucinating motion for absent objects.

\paragraph{Conditioning.}
The global condition vector fuses the diffusion timestep embedding with a CLIP~\cite{radford2021learning} ViT-B/32 text embedding:
\begin{equation}
    \mathbf{c} = \text{MLP}([\text{TimestepEmbed}(t); W_{\text{text}}\,\text{CLIP}(P)]).
    \label{eq:condition}
\end{equation}
Classifier-free guidance is enabled by randomly masking the text embedding with probability 0.1 during training.

\paragraph{Output projections.}
After the transformer, tokens are unflattened and projected back to raw dimensions through per-entity output layers. Crucially, the output predicts only the \emph{dynamic} components (position + velocity); BPS geometry is never predicted since it is static:
\begin{equation}
    \hat{\mathbf{v}} = [W_{\text{body}}^{\text{out}}\,\mathbf{h}_{\text{body}}; W_{\text{lh}}^{\text{out}}\,\mathbf{h}_{\text{lh}}; W_{\text{rh}}^{\text{out}}\,\mathbf{h}_{\text{rh}}; W_{\text{obj}}^{\text{out}}\,\mathbf{h}_{\text{obj}_1}; \ldots] \in \mathbb{R}^{534}.
    \label{eq:output}
\end{equation}

\paragraph{Flow matching objective.}
We train with Conditional Flow Matching~\cite{lipman2023flow}, which learns a velocity field $v_\theta$ transporting samples from a Gaussian prior $\mathbf{x}_0 \sim \mathcal{N}(0, \mathbf{I})$ to the data distribution:
\begin{equation}
    \mathcal{L}_{\text{FM}} = \mathbb{E}_{t, \mathbf{x}_0, \mathbf{x}_1}\left[\left\| v_\theta(\mathbf{x}_t, t, \mathbf{S}, \mathbf{c}, \mathbf{B}) - (\mathbf{x}_1 - \mathbf{x}_0) \right\|^2\right],
    \label{eq:fm_loss}
\end{equation}
where $\mathbf{x}_t = (1{-}t)\mathbf{x}_0 + t\,\mathbf{x}_1$ and $\mathbf{B}$ denotes the BPS embeddings. The loss is computed on the 534D dynamic features only.


\subsection{Block-wise Autoregressive Relay}
\label{sec:autoreg}

To generate seamless, long-horizon motions spanning multiple action phases, we employ a block-wise autoregressive strategy.

\paragraph{History-target decomposition.}
During training, each window of $T{=}200$ frames is decomposed into:
\begin{itemize}
    \item \textbf{Kinematic history} $\mathbf{X}_{\text{hist}}$: the first $K{=}20$ frames (2 seconds at 10~FPS) are kept \emph{noise-free}, providing a clean conditioning context with ground-truth positions and velocities.
    \item \textbf{Generation target} $\mathbf{X}_{\text{target}}$: the remaining $L{=}180$ frames receive flow matching noise and are the prediction target.
\end{itemize}

The history and noised target are concatenated along the temporal axis:
\begin{equation}
    \mathbf{x}_t^{\text{full}} = [\mathbf{X}_{\text{hist}}; \mathbf{x}_t^{\text{target}}] \in \mathbb{R}^{T \times 534},
    \label{eq:autoreg_concat}
\end{equation}
and the loss mask zeroes gradients for the $K$ history frames:
\begin{equation}
    \mathcal{L}_{\text{AR}} = \mathbb{E}\left[\left\| v_\theta(\mathbf{x}_t^{\text{full}}, t, \mathbf{S}, \mathbf{c}, \mathbf{B}) - \mathbf{u}_t \right\|^2 \cdot \mathbf{M}_{\text{target}}\right],
    \label{eq:autoreg_loss}
\end{equation}
where $\mathbf{M}_{\text{target}}$ is 0 for the first $K$ frames and 1 for the remaining $L$ frames.

\paragraph{Sliding window relay.}
For long-horizon generation ($T > 200$), CRR-Flow chains multiple generation windows:
\begin{enumerate}
    \item \textbf{Initialization}: $\mathbf{X}_{\text{hist}}$ is set to the rest pose with zero velocity for the first segment.
    \item \textbf{Generation}: The model generates $L{=}180$ frames conditioned on $\mathbf{X}_{\text{hist}}$ and the current state/text instruction.
    \item \textbf{Relay}: The \emph{last} $K$ frames of the generated segment become the new $\mathbf{X}_{\text{hist}}$, and the state planner advances to the next instruction.
\end{enumerate}
This ``handover'' mechanism preserves kinematic continuity: the hand's momentum at the end of a \textsc{Reach} phase flows naturally into the \textsc{Grasped} phase, and object positions are continuous across segment boundaries.


%% ============================================================
%% 4. EXPERIMENTS
%% ============================================================
\section{Experiments}
\label{sec:experiments}

\subsection{Datasets}
\label{sec:datasets}

\paragraph{OakInk2~\cite{zhan2024oakink2}.}
A bimanual hand-object manipulation dataset with full SMPL-X body annotations, MANO hand parameters, and multi-object 6DoF poses. We extract 963 super-segments from 627 complex sequences by merging adjacent primitives with temporal gaps $\leq 300$ frames. After boundary expansion ($\pm 60$ frames), 30$\rightarrow$10~FPS downsampling, and uniform subsampling to $T{=}200$ frames, each sample is a 534D dynamic feature vector (267D positions + 267D velocities) with 4 object slots. State labels $[T, 7]$ are computed via Schmitt trigger contact detection. Train/test split: 770/193 (80/20).

\paragraph{GRAB~\cite{taheri2020grab}.}
A single-hand grasping dataset with 1,335 sequences across 10 subjects and 51 objects. Raw SMPL-X parameters at 120~FPS are downsampled to 30~FPS, with ground-truth per-vertex contact labels on 10,475 SMPL-X vertices. Target length: 300 frames. Object slots 2--4 are padded. Subject-unseen split: train on s1--s8 (1,065), test on s9--s10 (270).

\paragraph{OMOMO.}
Full-body manipulation of large objects (furniture, appliances) with 4,294 sequences at 30~FPS from 17 subjects. Features are extracted from InterAct's~\cite{interact2024} unified 962D representation by selecting the compatible 534D subset (marker positions/velocities + object pose/velocity). 3,864 train / 430 test.

\paragraph{InterAct GRAB.}
InterAct's preprocessed version of GRAB (1,347 sequences) in canonical 962D format, converted to V4 534D. 1,077 train / 270 test.


\subsection{Implementation Details}
\label{sec:implementation}

\paragraph{Architecture.}
CRR-Flow uses a DiT backbone with latent dimension $D{=}512$, 8 transformer layers, 8 attention heads, and FFN dimension 2048. The spatio-temporal sequence length is $T {\times} 7$ tokens (e.g., $200 {\times} 7 = 1400$ for OakInk2). Total trainable parameters: 41.2M (excluding frozen CLIP ViT-B/32 at 151.3M).

CLIP text embeddings are pre-computed and cached, allowing the CLIP model to be excluded from the training graph for memory efficiency and DDP compatibility.

\paragraph{Training.}
We train with the Adam optimizer ($\text{lr}{=}10^{-4}$, weight decay $0.01$) for 200K steps with gradient clipping at 1.0. Batch size is 32 (single GPU) or $4 {\times} 64 = 256$ (4-GPU DDP). Flow matching uses independent coupling (ConditionalFlowMatcher with $\sigma{=}0$). The autoregressive variant uses $K{=}20$ history frames with masked loss.

For weight initialization, input projections use Xavier uniform, while output projections and the final layer use Xavier with gain $0.01$ (not zero, to prevent dead gradients from AdaLN-Zero gating).

\paragraph{Inference.}
Motion is generated via ODE integration of the learned velocity field using Euler's method with 10 steps from $t{=}0$ (noise) to $t{=}1$ (clean motion). For autoregressive generation, each window produces $L{=}180$ new frames, with the last $K{=}20$ frames relayed as history for the next window.

\paragraph{Visualization.}
Generated 77-marker positions are converted to SMPL-X mesh via LBFGS inverse fitting (two-stage: translation/orientation first, then full body). Object meshes are rendered with per-frame 6DoF transforms. Savitzky-Golay temporal smoothing (window $= 11$, polyorder $= 3$) is applied before fitting.


\subsection{Evaluation Metrics}
\label{sec:metrics}

We evaluate on three categories of metrics:

\paragraph{Generation quality.}
\textbf{FID}: Fr\'{e}chet Inception Distance on motion embeddings using InterAct's pre-trained encoder.
\textbf{R-Precision@k} ($k{=}1,2,3$): text-motion retrieval accuracy.
\textbf{Diversity}: pairwise distance variance in embedding space.

\paragraph{Reconstruction quality.}
\textbf{Marker MSE}: mean squared error of 77 marker positions vs.\ ground truth.
\textbf{Object Pose Error}: position (mm) and rotation (degrees) error.

\paragraph{Interaction quality (state-aware).}
\textbf{Contact F1}: precision/recall of hand-object contact at 5cm threshold, evaluated separately for each hand.
\textbf{State Transition Accuracy}: whether generated motion reflects state boundary timing (e.g., deceleration at \textsc{Reach}$\rightarrow$\textsc{Grasped}).
\textbf{Grasp Onset Error}: frame offset between generated and ground-truth first contact.

\subsection{Results}
\label{sec:results}

\paragraph{Training convergence.}
\Cref{tab:convergence} shows training loss across datasets. All models converge from an initial loss of ${\sim}2.0$ (expected for flow matching with normalized data) to well below 0.5 within 10K steps.

\begin{table}[t]
  \caption{Training loss convergence across datasets at various checkpoints. Loss is MSE on the 534D velocity field.}
  \label{tab:convergence}
  \begin{tabular}{@{}lcccccc@{}}
    \toprule
    Dataset & 500 & 2K & 5K & 10K & 50K \\
    \midrule
    OakInk2 (V3, 267D)  & 1.52 & 0.77 & 0.37 & 0.11 & 0.05 \\
    OakInk2 (V4, 534D)  & 1.58 & 0.85 & 0.36 & 0.27 & 0.15 \\
    OakInk2 (V4+AR)     & 1.56 & 0.86 & 0.37 & 0.27 & -- \\
    GRAB (V4, 30fps)    & 1.36 & 0.68 & 0.28 & 0.12 & -- \\
    InterAct GRAB        & 1.45 & 0.68 & -- & -- & -- \\
    OMOMO                & 1.49 & 0.48 & -- & -- & -- \\
    \bottomrule
  \end{tabular}
\end{table}

\paragraph{Generation quality at 10K steps.}
We evaluate early checkpoints (10K steps) on unseen test sequences using ground-truth state labels. \Cref{tab:gen_quality} compares V3 (position-only, 267D) and V4 (position+velocity+BPS, 534D).

\begin{table}[t]
  \caption{Marker MSE ($\times 10^{-2}$) on 4 OakInk2 test samples at 10K steps. V4 outperforms V3 on 3/4 samples. Lower is better.}
  \label{tab:gen_quality}
  \begin{tabular}{@{}lcccc@{}}
    \toprule
    & Sample 0 & Sample 1 & Sample 2 & Sample 3 \\
    \midrule
    V3 (267D, no BPS) & 0.79 & 0.94 & \textbf{4.75} & 1.11 \\
    V4 (534D + BPS)   & \textbf{0.98} & \textbf{0.56} & 1.97 & \textbf{0.85} \\
    \bottomrule
  \end{tabular}
\end{table}

V4's velocity features and BPS geometry provide stronger inductive bias for object interaction, particularly for samples with significant motion range (Sample 2: V3 collapses to 0.091 range vs.\ GT 0.188, while V4 maintains 0.155).

\paragraph{Ablation: autoregressive vs.\ non-autoregressive.}
At the same training step (4K), the autoregressive variant achieves comparable loss (0.375) to the non-autoregressive baseline (0.437), while enabling long-horizon relay generation. The $K{=}20$ history conditioning adds negligible computational cost (clean frames require no gradient computation).

\paragraph{Cross-dataset scaling.}
CRR-Flow V4 trains successfully across 4 datasets with different characteristics (bimanual vs.\ single-hand, small vs.\ large objects, 10~FPS vs.\ 30~FPS), demonstrating the generality of the entity-centric architecture and state conditioning framework.

%% Placeholder for full quantitative results table
%% \begin{table}[t]
%%   \caption{Full quantitative results (to be completed after 200K training).}
%%   \label{tab:full_results}
%% \end{table}
